\documentclass[11pt,a4paper]{article}
\usepackage[utf8]{inputenc}
\usepackage[T1]{fontenc}
\usepackage[spanish]{babel}
\usepackage{amsmath,amssymb,mathtools}
\usepackage{geometry}
\usepackage{enumitem}
\usepackage{array}
\usepackage{lmodern}
\usepackage{microtype}
\geometry{margin=2.4cm}
\setlength{\parindent}{0pt}
\setlength{\parskip}{0.45em}

\newcommand{\R}{\mathbb{R}}
\newcommand{\C}{\mathbb{C}}
\newcommand{\K}{\mathbb{K}}
\newcommand{\Ker}{\operatorname{Ker}}
\newcommand{\Ima}{\operatorname{Im}}
\newcommand{\rg}{\operatorname{rg}}
\newcommand{\tr}{\operatorname{tr}}
\newcommand{\Id}{\operatorname{Id}}
\newcommand{\diag}{\operatorname{diag}}
\newcommand{\Span}{\operatorname{L}}

\begin{document}

\begin{center}
{\LARGE\bfseries Álgebra Lineal}\\[0.4em]
{\Large Tema 3: Endomorfismos. Diagonalización}
\end{center}

\tableofcontents
\newpage

\section{Endomorfismos}

Un \textbf{endomorfismo} es una aplicación lineal:

\[
f:V\longrightarrow V
\]

de un espacio vectorial en sí mismo.

Si:

\[
\dim V=n
\]

la matriz asociada a $f$ en cualquier base es cuadrada de orden $n$.

Para un endomorfismo $f$ con matriz asociada $A$ son equivalentes:

\[
\Ker f=\{0\}
\]

\[
\dim(\Ima f)=n
\]

\[
\rg A=n
\]

\[
\det A\neq 0
\]

\[
f\text{ es inyectivo}
\iff
f\text{ es suprayectivo}
\iff
f\text{ es biyectivo}
\]

\section{Matrices semejantes}

Si $A$ y $M$ son matrices asociadas al mismo endomorfismo en bases diferentes, existe una matriz regular $P$ tal que:

\[
\boxed{
M=P^{-1}AP
}
\]

Dos matrices relacionadas de esta forma se denominan \textbf{semejantes}.

Las matrices semejantes tienen el mismo:

\[
\boxed{
\rg,\qquad \tr,\qquad \det
}
\]

y el mismo polinomio característico.

\section{Autovalores y autovectores}

Sea $A$ la matriz de un endomorfismo.

Un vector no nulo $v$ es un \textbf{autovector} asociado al escalar $\lambda$ si:

\[
\boxed{
Av=\lambda v
}
\]

El escalar $\lambda$ se denomina \textbf{autovalor}.

Equivalentemente:

\[
(A-\lambda I)v=0
\]

Para que exista una solución no trivial:

\[
\boxed{
\det(A-\lambda I)=0
}
\]

\subsection{Polinomio característico}

El polinomio:

\[
\boxed{
p(\lambda)=\det(A-\lambda I)
}
\]

se denomina polinomio característico.

Sus raíces son los autovalores del endomorfismo.

\subsection{Subespacio propio}

El subespacio propio correspondiente a $\lambda$ es:

\[
\boxed{
S(\lambda)=\Ker(A-\lambda I)
}
\]

Sus vectores no nulos son los autovectores asociados a $\lambda$.

\subsection{Propiedades}

Si $A$ es triangular o diagonal, sus autovalores son sus elementos diagonales.

Si $\lambda$ es autovalor de $A$, entonces:

\[
\lambda^k
\]

es autovalor de $A^k$,

\[
\alpha\lambda
\]

es autovalor de $\alpha A$,

\[
\lambda-k
\]

es autovalor de $A-kI$, y, si $A$ es regular,

\[
\frac1\lambda
\]

es autovalor de $A^{-1}$.

Además, $A$ y $A^t$ tienen los mismos autovalores.

Los autovectores asociados a autovalores distintos son linealmente independientes.

\section{Multiplicidad algebraica y geométrica}

Sea $\lambda$ un autovalor.

Su \textbf{multiplicidad algebraica} es la multiplicidad de $\lambda$ como raíz de:

\[
p(\lambda)
\]

y se denota $a_\lambda$.

Su \textbf{multiplicidad geométrica} es:

\[
\boxed{
g_\lambda=\dim S(\lambda)
=
\dim\Ker(A-\lambda I)
}
\]

Se verifica:

\[
\boxed{
1\leq g_\lambda\leq a_\lambda
}
\]

\section{Diagonalización}

Un endomorfismo es diagonalizable si existe una base:

\[
B=\{v_1,\ldots,v_n\}
\]

formada por autovectores.

En esa base:

\[
f(v_i)=\lambda_i v_i
\]

y:

\[
\boxed{
M_B(f)
=
D=
\diag(\lambda_1,\ldots,\lambda_n)
}
\]

Si $A$ es la matriz en la base canónica y $P$ tiene por columnas los autovectores de la base diagonalizadora:

\[
\boxed{
D=P^{-1}AP
}
\]

o equivalentemente:

\[
\boxed{
A=PDP^{-1}
}
\]

\subsection{Criterio de diagonalización}

Una matriz $A$ de orden $n$ es diagonalizable si y sólo si:

\[
\boxed{
\sum_\lambda a_\lambda=n
}
\]

y para cada autovalor:

\[
\boxed{
g_\lambda=a_\lambda
}
\]

En particular, si $A$ tiene $n$ autovalores distintos, entonces es diagonalizable.

\section{Diagonalización ortogonal}

Una matriz $P$ es ortogonal si:

\[
P^{-1}=P^t
\]

Una matriz real $A$ es ortogonalmente diagonalizable si existe una matriz ortogonal $P$ tal que:

\[
\boxed{
D=P^tAP
}
\]

El teorema espectral establece:

\[
\boxed{
A\text{ real y simétrica}
\iff
A\text{ es ortogonalmente diagonalizable}
}
\]

Los autovectores de una matriz simétrica asociados a autovalores distintos son ortogonales.

\section{Subespacios invariantes}

Un subespacio $S\leq V$ es \textbf{invariante} por $f$ si:

\[
\boxed{
f(S)\subseteq S
}
\]

Todo subespacio propio:

\[
S(\lambda)=\Ker(f-\lambda I)
\]

es invariante.

Para estudiar cadenas de Jordan se utilizan los subespacios:

\[
\boxed{
E_k(\lambda)=\Ker(f-\lambda I)^k
}
\]

que verifican:

\[
E_1(\lambda)\subseteq
E_2(\lambda)\subseteq
\cdots
\]

hasta estabilizar en el subespacio característico asociado a $\lambda$.

\section{Teorema de Cayley--Hamilton}

Toda matriz cuadrada satisface su propio polinomio característico.

Si:

\[
p(\lambda)
=
a_0+a_1\lambda+\cdots+a_n\lambda^n
\]

entonces:

\[
\boxed{
p(A)
=
a_0I+a_1A+\cdots+a_nA^n
=
0
}
\]

Este resultado permite reducir potencias altas de matrices.

Si $\det A\neq0$, también permite expresar $A^{-1}$ como combinación lineal de:

\[
I,A,\ldots,A^{n-1}
\]

\section{Forma canónica de Jordan}

Un bloque de Jordan de orden $r$ asociado a $\lambda$ es:

\[
J_r(\lambda)
=
\begin{pmatrix}
\lambda&0&\cdots&0\\
1&\lambda&\ddots&\vdots\\
0&1&\ddots&0\\
\vdots&\ddots&\ddots&\lambda
\end{pmatrix}
\]

La forma canónica de Jordan es una matriz diagonal por bloques:

\[
\boxed{
J=
\diag
\left(
J_{r_1}(\lambda_1),
\ldots,
J_{r_s}(\lambda_s)
\right)
}
\]

tal que:

\[
\boxed{
J=P^{-1}AP
}
\]

El número de bloques asociados a $\lambda$ coincide con:

\[
\boxed{
\dim\Ker(A-\lambda I)
}
\]

y el tamaño del mayor bloque es el menor entero $k$ para el que:

\[
\Ker(A-\lambda I)^k
\]

alcanza el subespacio característico.

\subsection{Cadenas de Jordan}

Si:

\[
v_1\in
\Ker(A-\lambda I)^r
\setminus
\Ker(A-\lambda I)^{r-1}
\]

se construye la cadena:

\[
v_1,
\quad
v_2=(A-\lambda I)v_1,
\quad
\ldots,
\quad
v_r=(A-\lambda I)^{r-1}v_1
\]

que genera un bloque de Jordan de orden $r$.

\section{Forma real de Jordan}

Si $A$ es real y tiene un autovalor complejo:

\[
\lambda=a+bi
\]

también aparece:

\[
\overline{\lambda}=a-bi
\]

La pareja se representa en forma real mediante el bloque:

\[
\boxed{
C(\lambda)
=
\begin{pmatrix}
a&b\\
-b&a
\end{pmatrix}
}
\]

Los bloques reales de Jordan correspondientes se construyen a partir de estos bloques $2\times2$.

\section{Exponencial de una matriz}

La exponencial de una matriz se define como:

\[
\boxed{
e^A
=
I+A+\frac{A^2}{2!}
+\frac{A^3}{3!}
+\cdots
}
\]

Si:

\[
A=PJP^{-1}
\]

entonces:

\[
\boxed{
e^A
=
Pe^JP^{-1}
}
\]

Para un bloque:

\[
J_r(\lambda)
=
\lambda I+N
\]

con $N^r=0$:

\[
\boxed{
e^{J_r(\lambda)}
=
e^\lambda
\left(
I+N+\frac{N^2}{2!}
+\cdots+
\frac{N^{r-1}}{(r-1)!}
\right)
}
\]

En particular:

\[
e^{\lambda I}=e^\lambda I
\]

\section{Factorización LU y LDU}

La factorización LU expresa una matriz como:

\[
\boxed{
A=LU
}
\]

donde $L$ es triangular inferior y $U$ triangular superior.

La factorización LDU separa además los elementos diagonales:

\[
\boxed{
A=LDU
}
\]

donde:

\[
L
\]

es triangular inferior invertible,

\[
D
\]

es diagonal invertible, y:

\[
U
\]

es triangular superior invertible.

La idea consiste en obtener primero una factorización LU y extraer la diagonal de $U$ para formar $D$.

\end{document}
